\begin{lemma}{Injektivit�t linearer Abbildungen zwischen norm. R�umen}
             {Injektivit�tLinearerAbbildungenZwischenNormiertenR�umen}
Seien $(E, \norm\cdot_E), (F, \norm\cdot_F)$ normierte $\K$--Vektorr�ume und
$T \in \Hom(E, F)$ mit der Eigenschaft
$\inf\menge{\norm{Tx}_F}{x \in \kran E 0 1} \neq 0$.\\
Dann ist $T$ injektiv.
\end{lemma}
\begin{proof}
Sei $T$ nicht injektiv. Dann gibt es ein $x\in E\setminus\meng 0$ mit $Tx=0$.\\
Da $\norm x _E > 0$ ist $x' := \frac 1 {\norm x_E}x$ wohldefiniert.\\
Es ist also $\norm {x'}_E = 1$ und 
$\norm{Tx'}_F = \frac 1 {\norm x _E}\norm{Tx}_F = 0$, da $Tx = 0$.\\
Wegen 
$\menge{\norm{Ta}_F}{a \in \kran E 0 1} \subseteq [0,\infty[$,
ist 0 eine untere Schranke dieser Menge und nach obigem auch darin
enthalten, also gleich ihrem Infimum.
\end{proof}